\subsection{实值函数的导数}

当我们处理实变量的实值函数时，前面的定义和命题可以在若干方面加以扩充。

首先，若 \(f\) 是这样一个函数：定义在区间 \(I \subset \mathbf{R}\) 上、且相对于 \(I\) 在一点 \(x_0 \in I\) 连续，则当 \(x\) 趋于 \(x_0\) 且保持属于 \(I\) 并 \(\neq x_0\) 时，\(\frac{f(x) - f(x_0)}{x - x_0}\) 可能以 \(+\infty\) 或 \(-\infty\) 为极限；这时称 \(f\) 在 \(x_0\) 可导并在该点具有导数 \(+\infty\) （相应地，\(-\infty\) ）；若函数 \(f\) 具有导数 \(f'(x)\) （有限或无穷）在每一点 \(x\) （属于 \(I\) ）处，则函数 \(f'\) （取值于 \(\overline{\mathbf{R}}\) ）仍称为 \(f\) 的导函数（或简称导数）。右导数与左导数的定义同样推广。

例。在点 x = 0 处，函数 \(x^{1/3}\) （\(x^{3}\) 的反函数，后者是 R 到自身上的同胚）具有导数，等于 \(+\infty\) ；在 x = 0 处函数 \(|x|^{1/3}\) 具有右导数 \(+\infty\) 与左导数 \(-\infty\) 。

可导实函数的和、积以及可导函数的逆（命题 1、3 与 4）的导数公式，以及（实值）复合函数的导数（命题 5）公式，当所涉及的导数为无穷时仍然成立，只要这些公式中出现的所有表达式都有意义（Gen.~Top., IV, p.~345--346）。事实上，若在命题 6 中假设 f 在 I 上严格递增（相应地，严格递减）且连续，且 \(f'(x_{0}) = 0\) ，则反函数 g 具有导数，等于 \(+\infty\) （相应地，\(-\infty\) ）（在点 \(y_{0} = f(x_{0})\) 处）；若 \(f'(x_{0}) = +\infty\) （相应地，\(-\infty\) ），则 g 具有导数 0。对右导数与左导数有类似的结果，我们把它留给读者。

设 C 是有限实函数 f 的图像即表示曲线，它是平面 \(R^{2}\) 中由点 \((x, f(x))\) （x 跑遍 f 的定义集）所成的子集。若函数 f 在点 \(x_{0} \in I\) 具有有限的右导数，则以 C 的点 \(\mathbf{M}_{x_{0}} = (x_{0}, f(x_{0}))\) 为原点、方向数为 \((1, f_{d}'(x_{0}))\) 的半直线称为曲线 C 在点 \(M_{x_{0}}\) 处的右半切线；当 \(f_{d}'(x_{0}) = +\infty\) （相应地，\(f_{d}'(x_{0}) = -\infty\) ）时，我们对以 \(M_{x_{0}}\) 为原点、方向数为 \((0, 1)\) （相应地，\((0, -1)\) ）的半直线沿用同样的术语。同样地，可定义 \(M_{x_{0}}\) 处的左半切线（当 \(f_{g}'(x_{0})\) 存在时）。利用这些定义可以迅速验证：右（相应地，左）半切线与横轴的夹角，是该轴与以 \(M_{x_{0}}\) 为原点且通过 C 的点 \(\mathbf{M}_{x} = (x, f(x))\) 的半直线所夹之角当 x 趋于 \(x_{0}\) 且保持 \textgreater{} \(x_{0}\) （相应地，\textless{} \(x_{0}\) ）时的极限。

也可以说，右（相应地，左）半切线是以 \(M_{x_{0}}\) 为原点且通过 \(M_{x}\) 的半直线的极限，这时对具有同一原点的半直线所成的集赋予商空间 \(C^{*}/R_{+}^{*}\) 的拓扑（Gen.~Top., VIII, p.~107）。

若两条半切线在 C 的一点 \(M_{x_{0}}\) 处都存在，则仅当 f 在点 \(x_{0}\) （设为 I 的内点）具有导数（有限或无穷）时它们才方向相反；仅当 \(f_{d}^{\prime}(x_{0})\) 与 \(f_{g}^{\prime}(x_{0})\) 均为无穷且符号相反时它们才重合。在这两种情形，我们称包含这两条半切线的直线为 C 在点 \(M_{x_{0}}\) 的切线。

当 \(M_{x_{0}}\) 处的切线存在时，它是通过 \(M_{x_{0}}\) 与 \(M_{x}\) 的直线当 x 趋于 \(x_{0}\) 且保持 \(\neq x_{0}\) 时的极限，过一定点的直线所成之集上的拓扑是商空间 \(C^{*}/R^{*}\) 的拓扑（Gen.~Top., VIII, p.~114）。

图像的切线与半切线概念是一般概念的特例，后者将在本丛书专讲微分流形的部分中定义。

定义 4. 设 \(f\) 是定义在一个子集 \(A\) （拓扑空间 \(E\) 的子集）上的实函数；称 f 在一点 \(x_0 \in A\) 处（相对于 \(A\) ）具有相对极大值（相应地，严格相对极大值、相对极小值、严格相对极小值），如果存在一个邻域 \(V\) （\(x_0\) 在 \(E\) 中的邻域），使得在每个点 \(x \in V \cap A\) （异于 \(x_0\) ）处都有 \(f(x) \leqslant f(x_0)\) （相应地，\(f(x) < f(x_0)\) 、\(f(x) \geqslant f(x_0)\) 、\(f(x) > f(x_0)\) ）。

显然，若 f 在 A 的一点处达到它在 A 上的上确界（相应地，下确界），则它在该点相对于 A 具有相对极大值（相应地，相对极小值）；反之当然不成立。

注意，若 \(B \subset A\) ，且 f （例如）在一点 \(x_{0} \in B\) 处相对于 B 具有相对极大值，则 f 在该点未必相对于 A 具有相对极大值。

命题 7. 设 \(f\) 是定义在区间 \(\mathbf{I} \subset \mathbf{R}\) 上的有限实函数。若 \(f\) 在一个内点 \(x_0\) （\(\mathbf{I}\) 的内点）处具有相对极大值（相应地，相对极小值），且在该点既有右导数又有左导数，则有 \(f_d'(x_0) \leqslant 0\) 及 \(f_g'(x_0) \geqslant 0\) （相应地，\(f_d'(x_0) \geqslant 0\) 及 \(f_g'(x_0) \leqslant 0\) ）；特别地，若 \(f\) 在点 \(x_0\) 可导，则 \(f'(x_0) = 0\) 。

命题的证明由定义直接得出。

我们还可以说，若 f 在 I 的一个内点 \(x_{0}\) 处可导且具有相对极大值或相对极小值，则其图像的切线平行于横轴。反之不成立，函数 \(x^{3}\) 的例子表明了这一点：它在点 x = 0 处导数为零，但在该点既无相对极大值亦无相对极小值。

